% Generated draft: requires mathematical and source review.
\documentclass[11pt]{article}
\usepackage[margin=25mm]{geometry}
\usepackage{fontspec}
\setmainfont{Latin Modern Roman}
\usepackage{amsmath,amssymb,mathtools,unicode-math}
\setmathfont{Latin Modern Math}
\usepackage{enumitem,needspace,xurl,hyperref}
\setlength{\parindent}{0pt}
\setlength{\parskip}{5pt}
\setlength{\emergencystretch}{3em}
\newcommand{\N}{\mathbb{N}}
\newcommand{\Z}{\mathbb{Z}}
\newcommand{\Q}{\mathbb{Q}}
\newcommand{\R}{\mathbb{R}}
\newcommand{\C}{\mathbb{C}}
\newcommand{\F}{\mathbb{F}}
\newcommand{\A}{\mathbb{A}}
\newcommand{\PP}{\mathbb{P}}
\newcommand{\E}{\mathbb{E}}
\newcommand{\eps}{\varepsilon}
\newcommand{\dd}{\,\mathrm{d}}
\begin{document}
\section*{UnsolvedMath ID: 30006486}
\textbf{Source problem code:} catalogue entry
\subsection{Definitions and mathematical statement}
A set-theoretic solution of the Yang--Baxter equation consists of a set \(X\) and a map \(r:X\times X\to X\times X\) satisfying
\[
(r\times\operatorname{id}_X)(\operatorname{id}_X\times r)(r\times\operatorname{id}_X)
=
(\operatorname{id}_X\times r)(r\times\operatorname{id}_X)(\operatorname{id}_X\times r)
\quad\text{on }X^3.
\]
Write \(r(x,y)=(\sigma_x(y),\tau_y(x))\). Nondegeneracy means that every \(\sigma_x\) and every \(\tau_y\) is a permutation of \(X\); bijectivity of \(r\) is a separate condition. In the usual nonempty, bijective, nondegenerate setting, indecomposability means that the subgroup of \(\operatorname{Sym}(X)\) generated by all \(\sigma_x\) and \(\tau_x\) acts transitively on \(X\). The supplied excerpt does not specify whether this setting is intended, or whether finiteness or involutivity, \(r^2=\operatorname{id}_{X\times X}\), is required. These conditions must therefore remain unresolved qualifications.

The supplied record contains the following research questions.
\begin{enumerate}
\item Is there a systematic construction that recomposes, or stitches together, indecomposable set-theoretic solutions into new solutions? The permitted input families, compatibility data, meaning of recomposition, and required properties of the output are not specified. Thus no particular construction or universal existence assertion is determined by the excerpt.
\item Can solutions be classified up to cabling, or can invariants distinguish cabling-equivalence classes? The cleaned statement additionally asks for computable invariants. Formally, after specifying an intended class \(\mathcal C\) of solutions and a cabling equivalence relation \(\sim_{\mathrm{cab}}\), an invariant with values in a set \(S\) would be a map \(I:\mathcal C\to S\) satisfying
\[
A\sim_{\mathrm{cab}}B\ \Longrightarrow\ I(A)=I(B).
\]
A complete invariant would additionally satisfy the converse. The source does not define cabling, the resulting equivalence relation, or whether complete distinction of all classes is requested. Computability also requires an effective presentation of the solutions and invariant values, which is not supplied. It is unclear whether the classification concerns all intended solutions or only indecomposable ones.
\item The original excerpt also asks, under ``Structure algebras of set-theoretic solutions'': what can be said about the integral unit group of the multiplicative group of \(B\)? This question is absent from the cleaned statement but is retained here as a source qualification. Neither \(B\) nor its construction is defined. If the intended object is the unit group of the integral group ring of a group \(G\), its notation would be \((\mathbb Z[G])^\times\); the excerpt does not establish this interpretation or specify whether a subgroup of units is intended. Consequently this part cannot be made into a definite mathematical statement from the supplied evidence.
\end{enumerate}
These are statement-only research targets with unresolved hypotheses and operations, rather than a uniquely specified conjecture.
\subsection{Short English statement}
The record asks how to assemble indecomposable Yang--Baxter solutions and how to classify solutions using cabling or computable invariants. Its original excerpt also includes a question about an integral unit group associated with an undefined object \(B\). The intended solution class, assembly operation, cabling relation, and algebraic object require clarification before a precise conjecture can be extracted.
\subsection{Sources}
\begin{itemize}
\item Local UnsolvedMath registry, \texttt{data/unsolvedmath/problems.json}, permanent ID \(30006486\), source problem code \texttt{catalogue entry}, title ``Recomposition and Cabling of Set-Theoretic Yang--Baxter Solutions.'' No individual proposer is supplied; the proposed year is recorded as 2025.
\item Supplied source citation: \emph{Mini-Workshop: The Yang--Baxter Equation and Representations of Braid Groups} (2026), Oberwolfach Reports, pp. 2633--2672, DOI \url{https://doi.org/10.4171/owr/2025/49}. The literature-source label is ``Oberwolfach Report 49/2025,'' and the extracted passage includes page 2637. The difference between the citation year 2026 and the report designation 2025 is unresolved in the supplied record.
\item Supplied literature source: \emph{Cabling of non-involutive solutions}, \url{https://arxiv.org/abs/2410.23821}. No authors or specific results from this reference are supplied.
\item Literature assessment dated exactly \texttt{2026-08-22}: supplied dated curation, not independently verified here. It labels the problem open and reports that recomposition, integral-unit-group, and cabling-classification questions remain open, with partial constructions and criteria but no requested general classification. The background explicitly calls for independent verification of time-sensitive status and cited claims.
\item The registry assigns the statement status \texttt{corrected\_verified} and reports that the official report was checked and its notation repaired. These are supplied verification claims, not independent verification by this draft. The record attributes the questions to Oberwolfach workshop open-problem discussions.
\end{itemize}
\end{document}
